\documentclass[a4paper, serif, 10pt]{moderncv}

%% moderncv
% color and style
\moderncvcolor{black}
\moderncvstyle{banking}

% renew moderncv command to adapt for CJK colon
\renewcommand*{\cvitem}[3][.25em]{%
  \ifstrempty{#2}{}{\hintstyle{#2}：}{#3}%
  \par\addvspace{#1}}

\renewcommand*{\cvitemwithcomment}[4][.25em]{%
  \savebox{\cvitemwithcommentmainbox}{\ifstrempty{#2}{}{\hintstyle{#2}：}#3}%
  \setlength{\cvitemwithcommentmainlength}{\widthof{\usebox{\cvitemwithcommentmainbox}}}%
  \setlength{\cvitemwithcommentcommentlength}{\maincolumnwidth-\separatorcolumnwidth-\cvitemwithcommentmainlength}%
  \begin{minipage}[t]{\cvitemwithcommentmainlength}\usebox{\cvitemwithcommentmainbox}\end{minipage}%
  \hfill% fill of \separatorcolumnwidth
  \begin{minipage}[t]{\cvitemwithcommentcommentlength}\raggedleft\small\itshape#4\end{minipage}%
  \par\addvspace{#1}}

%% page layout/margins
\usepackage[top=2.5cm, bottom=2.5cm, left=1.5cm, right=1.5cm]{geometry}

\nopagenumbers{}

%% Babel config for English language
\usepackage[english]{babel}

%% fontspec
\usepackage{fontspec}

\IfFontExistsTF{Linux Libertine}{
  \setmainfont[Ligatures={TeX, Common}, Numbers=Lining, ItalicFont=Linux Libertine]{Linux Libertine}
}{}
\IfFontExistsTF{Linux Libertine O}{
  \setmainfont[Ligatures={TeX, Common}, Numbers=Lining, ItalicFont=Linux Libertine O]{Linux Libertine O}
}{}

%% CTeX
% CJK support, used to show CJK characters in the resume
%
% - heading=false: disable ctex heading
% - punct=kaiming: use kaiming punctuations styles for CJK
% - scheme=plain: use plain scheme, do not override \normalsize font size
% - space=auto: space settings for CJK characters
%
% CJK fonts are resolved per locale language, e.g. Noto Serif/Sans CJK JP
% for Japanese resumes, with fallback to the next region if not installed
%
% ref:
% - http://ctan.mirrorcatalogs.com/language/chinese/ctex/ctex.pdf
\usepackage[UTF8, heading=false, punct=kaiming, scheme=plain, space=auto]{ctex}

\IfFontExistsTF{Noto Serif CJK JP}{
  \setCJKmainfont{Noto Serif CJK JP}
}{}
\IfFontExistsTF{Noto Sans CJK JP}{
  \setCJKsansfont{Noto Sans CJK JP}
}{}

%% line spacing
\usepackage{setspace}
\setstretch{1.125}

%% URL styling - use normal text instead of monospace
\urlstyle{same}

%% Patch \httplink and \httpslink to handle full URLs
% The original moderncv macros always prepend http:// or https:// to URLs.
% This causes issues when the URL already has a protocol (e.g., https://example.com)
% resulting in malformed URLs like https://https://example.com.
% This patch checks if the URL already contains "://" and uses it directly if so.
% Note: We use \str_set:Nx (x-expansion) to expand macros like \@homepage first.
\ExplSyntaxOn
\str_new:N \g_yamlresume_url_str

\RenewDocumentCommand{\httpslink}{O{}m}{
  \str_set:Nx \g_yamlresume_url_str {#2}
  \str_if_in:NnTF \g_yamlresume_url_str {://}
    { \url{#2} }
    { \url{https://#2} }
}

\RenewDocumentCommand{\httplink}{O{}m}{
  \str_set:Nx \g_yamlresume_url_str {#2}
  \str_if_in:NnTF \g_yamlresume_url_str {://}
    { \url{#2} }
    { \url{http://#2} }
}
\ExplSyntaxOff

\name{高橋 健司}{}
\title{財務分析 / FP\&A / 企業価値評価}
\phone[mobile]{+81 90-1234-5678}
\email{kenji.takahashi.finance@example.jp}
\homepage{https://www.linkedin.com/in/kenji-takahashi-finance}

\address{日本、東京都、千代田区、東京都千代田区丸の内1-2-3、100-0005}{}{}

\extrainfo{{\small \faLinkedin}\ \href{https://www.linkedin.com/in/kenji-takahashi-finance}{@kenji-takahashi-finance} {} {} {} • {} {} {} 
{\small \faTwitter}\ \href{https://twitter.com/kenji\_fpna}{@kenji\_fpna}}

\begin{document}

\maketitle

\section{基本情報}

\cvline{}{\begin{itemize}
\item 企業価値評価とFP\&Aを専門とする財務アナリスト。7年以上の実務経験を通じ、予算策定、業績管理、M\&Aデューデリジェンス、経営層向け報告資料の作成に従事。
\item 財務モデリングとデータ分析を組み合わせ、経営判断に資する示唆を提供。Excel、SQL、Power BI、Pythonを活用した分析基盤の改善に強み。
\item 銀行および事業会社での経験を活かし、会計・財務・事業部門の橋渡しを推進。日本語と英語でのコミュニケーションが可能。
\end{itemize}}
\section{学歴}

\cventry{2014年4月～2018年3月}
        {学士、商学部金融学科、成績：3.7/4.0}
        {一橋大学}
        {\url{https://www.hit-u.ac.jp/}}
        {}
        {\begin{itemize}
\item 財務会計、コーポレートファイナンス、統計学を中心に履修
\item ゼミナールで企業価値評価と資本市場の研究に従事
\item グループワークを通じて財務分析とプレゼンテーション能力を養成
\end{itemize}
\textbf{コース}：財務会計、管理会計、コーポレートファイナンス、投資理論、統計学、計量経済学、企業価値評価、資本市場論}
\section{職歴}

\cventry{2021年4月～現在}
        {シニア・ファイナンシャル・アナリスト}
        {みらいキャピタル株式会社}
        {}
        {}
        {\begin{itemize}
\item 中期経営計画の策定支援として、部門別予算と5カ年財務モデルを構築し、経営会議向けのシナリオ分析を実施
\item M\&A案件において対象企業の財務デューデリジェンス、バリュエーション、PMI初期計画の策定を担当
\item 月次・四半期の業績管理を高度化し、KPIダッシュボードと差異分析レポートを経営層へ提供
\item 事業部門と連携して収益改善施策の定量評価を行い、原価低減と価格最適化の意思決定を支援
\item 若手アナリスト2名の育成を担当し、財務モデリングの品質基準とレビュー体制を整備
\end{itemize}
\textbf{キーワード}：財務モデリング、M\&A、予算管理、FP\&A、経営報告}

\cventry{2019年1月～2021年3月}
        {ファイナンシャル・アナリスト}
        {東都銀行}
        {}
        {}
        {\begin{itemize}
\item 法人融資案件の信用分析、財務諸表分析、返済能力評価を担当
\item 業種別のリスク要因を整理し、与信判断に必要な分析レポートを作成
\item 融資先の月次モニタリングを設計し、早期警戒指標の可視化を推進
\item Excel VBAとSQLを用いた集計業務の自動化により、月次作業時間を約40\%削減
\end{itemize}
\textbf{キーワード}：信用分析、融資審査、Excel、SQL、Power BI}

\cventry{2018年4月～2018年12月}
        {アナリスト}
        {サクラ会計事務所}
        {}
        {}
        {\begin{itemize}
\item 上場企業および中小企業の財務諸表分析、監査支援、税務申告補助を担当
\item 経理担当者へのヒアリングを通じて、改善提案と内部統制の整備を支援
\item 決算整理仕訳の確認と分析資料の作成を実施
\end{itemize}
\textbf{キーワード}：財務諸表分析、監査支援、IFRS、税務}
\section{言語}

\cvline{日本語}{ネイティブ・母国語レベル \hfill \textbf{キーワード}：日本語、ビジネス文書}
\cvline{英語}{完全な実務レベル \hfill \textbf{キーワード}：TOEIC 920、英語財務報告、海外ステークホルダー対応}
\section{スキル}

\cvline{財務分析}{エキスパート \hfill \textbf{キーワード}：財務モデリング、DCF法、企業価値評価、予算策定、差異分析}
\cvline{データ分析}{上級 \hfill \textbf{キーワード}：Excel、SQL、Power BI、Python、Tableau}
\cvline{会計・報告}{上級 \hfill \textbf{キーワード}：IFRS、JGAAP、連結決算、管理会計}
\cvline{コミュニケーション}{上級 \hfill \textbf{キーワード}：経営層向け報告資料作成、ステークホルダー調整、英語プレゼンテーション、チーム育成}
\section{受賞・表彰}

\cventry{2023年12月}
        {みらいキャピタル株式会社}
        {優秀プロジェクト賞}
        {}
        {}
        {M\&Aデューデリジェンスと買収後統合計画の定量分析が高く評価され、経営企画部門より表彰。}
\section{資格・免状}

\cventry{2022年2月}
        {CFA協会}
        {CFAレベル2合格}
        {\url{https://www.cfainstitute.org/}}
        {}
        {}

\cventry{2017年11月}
        {日本商工会議所}
        {日商簿記1級}
        {}
        {}
        {}

\cventry{2020年6月}
        {コーポレート・ファイナンス・インスティテュート}
        {FMVA（財務モデリング・バリュエーション・アナリスト）}
        {\url{https://corporatefinanceinstitute.com/}}
        {}
        {}
\section{出版・著作}

\cventry{2023年10月}
        {国内製薬業界のM\&A動向と企業価値評価}
        {みらいキャピタル調査部}
        {\url{https://www.example.com/research/pharma-ma-valuation}}
        {}
        {\begin{itemize}
\item 国内製薬企業のM\&A事例を対象に、シナジー評価とバリュエーション手法を比較分析
\item 創薬パイプラインの価値評価におけるリスク調整アプローチを整理
\item 経営企画・事業開発部門向けに、案件検討時の実務示唆を提供
\end{itemize}}
\section{プロジェクト}

\cventry{2023年1月～2023年6月}
        {経営層と事業部門向けに、予算・実績・予測を一元管理するFP\&Aダッシュボードを再構築}
        {FP\&Aダッシュボード刷新プロジェクト}
        {\url{https://www.example.com/projects/fpa-dashboard}}
        {}
        {\begin{itemize}
\item 部門別の予算消化率、売上・利益の差異要因、キャッシュフロー予測を可視化
\item SQLによるデータ整備とPower BIによるレポート自動更新を設計
\item 経営会議の意思決定に必要な分析リードタイムを従来比で約50\%短縮
\end{itemize}
\textbf{キーワード}：FP\&A、Power BI、SQL、予算管理、可視化}
\section{興味・関心}

\cvline{読書}{企業価値評価、経済史、英語ビジネス書}
\cvline{スポーツ}{テニス、ランニング、登山}
\cvline{金融教育}{金融リテラシー、コミュニティ活動}
\section{ボランティア活動}

\cventry{2022年9月～現在}
        {金融リテラシー講師}
        {NPO法人金融教育推進ネットワーク}
        {\url{https://www.example.com/volunteer/financial-literacy}}
        {}
        {\begin{itemize}
\item 高校生・大学生向けに、家計管理、資産形成、投資の基礎をテーマとしたワークショップを実施
\item 教材作成とグループディスカッションのファシリテーションを担当
\item 参加者アンケートの分析をもとに、講座内容の改善を継続
\end{itemize}}
    
\end{document}